\section{Methods and modelled assumptions}
\label{app:methods}

This appendix records the measurement-model conventions, the ephemeris handling, the
constellation model, and the numerical method behind the results. All are consistent
with the vendored provenance notice (Section~\ref{sec:availability}).

\subsection{Measurement model and projectors}
\label{app:model}

The geometry matrix $\Gmat$ of \eqref{eq:measmodel} uses the sign convention that the
position partial of the $i$-th pseudorange residual is the negative user-to-satellite
LOS unit vector $-\los_i$ and the clock partial is unity, so that a satellite-position
error $\ds_i$ enters the residual as $+\los_i\cdot\ds_i$ (an ephemeris error appears as
a range error along the sightline). The range and parity projectors
\eqref{eq:projectors} are formed from the thin QR factorization of $\Gmat$: with
$\Gmat=\bm Q\bm R$, $\bm Q\in\R^{n\times4}$ orthonormal, $\Ppar=\bm Q\bm Q^\top$ and
$\Pperp=\bm I_n-\bm Q\bm Q^\top$, which is numerically preferable to forming
$(\Gmat^\top\Gmat)^{-1}$ explicitly and is the route both the engine and the NumPy oracle
take. A parity matrix $\Smat$ (any orthonormal basis of $\rng(\Gmat)^\perp$, e.g.\ from
the full QR) gives the same $\Pperp=\Smat^\top\Smat$; the split quantities
\eqref{eq:blinddx}--\eqref{eq:detnorm} are basis-independent and we report them as such.

\subsection{Ephemeris provenance and access}
\label{app:provenance}

Three independent authoritative ephemeris families supply the real common-shift
magnitude, read through the IMCCE \code{calceph} library so that all three are accessed
by one interface in a single frame and timescale:
\begin{itemize}
  \item \textbf{DE440} (JPL/NASA), kernel \code{de440s.bsp} \citep{park2021de440}, from
  \url{https://naif.jpl.nasa.gov/pub/naif/generic_kernels/spk/planets/de440s.bsp};
  \item \textbf{INPOP21a} (IMCCE / Observatoire de Paris)
  \citep{fienga2021inpop21a}, from \url{https://ftp.imcce.fr/pub/ephem/planets/inpop21a/};
  \item \textbf{EPM2021} (IAA~RAS), kernel \code{epm2021.bsp}
  \citep{pitjeva2021epm2021}, from \url{https://ftp.iaaras.ru/pub/epm/EPM2021/SPICE/}.
\end{itemize}
Following the vendoring convention of the companion papers \citep{baweja2026datum,%
baweja2026interop}, the kernels are \emph{not} redistributed; only the sampled derived
positions and the computed reference split are vendored. Published scientific ephemeris
positions are factual constants and are not copyrightable; the derived sample points are
cited with full attribution to JPL, IMCCE, and IAA~RAS.

\subsection{Frame, timescale, and sampling}
\label{app:sampling}

The geocentric Moon position (NAIF body $301$ with respect to $399$) is sampled from
each family over 2024-01-01 to 2025-12-31 in Barycentric Dynamical Time (TDB) at two-day
cadence ($366$ epochs), in the ICRF/J2000 frame, so the three families are directly
comparable without a frame or timescale conversion. For each provider pair and epoch, the
per-epoch position difference is the common frame/ephemeris shift $\ds$ applied to the
constellation. Positions are rounded to six decimals (micrometre resolution) before the
fit in both the engine and the oracle, so the comparison tests the two \emph{algorithms}
at byte-identical inputs and is not sensitive to sampling-level differences. The record
count is $732$ ($2$ provider pairs $\times366$ epochs); Table~\ref{tab:anchor} reports
the per-pair medians.

\subsection{The modelled constellation}
\label{app:constellation}

The geometry matrix $\Gmat$ is built from a modelled lunar navigation constellation
(satellite LOS unit vectors to a representative user), used to project the common shift
$\ds$ into measurement space. This constellation is \Modelled{}: it stands in for a real
provider's beacon geometry, which has not flown. The blind-subspace membership of the
common shift (Lemma~\ref{lem:rigid}) is \emph{independent} of the constellation geometry
and of the user position---it holds for any $\Gmat$ of the form \eqref{eq:measmodel}---so
the modelled geometry affects only the partition of the fixed $\approx\!2\,$m magnitude
between the horizontal and vertical position components, not the fact of blindness. The
horizontal/vertical split feeding $\sigma_H$ in \eqref{eq:sigmaH} is modelled as
isotropic ($\Sig_{\mathrm{cm}}=\tfrac{1}{3}\sigma_{\mathrm{cm}}^2\bm I_3$); a specific
constellation and user geometry would replace this isotropic assumption with the actual
position-domain covariance, shifting $\sigma_H$ but not the mechanism.

\subsection{The CMPL and envelope}
\label{app:cmpl}

The common-mode one-sigma $\sigma_{\mathrm{cm}}=\SI{2}{m}$ is fixed to the real
inter-ephemeris floor of Section~\ref{sec:anchor} and \citep{baweja2026interop}. The
horizontal one-sigma follows from the isotropic assumption as
$\sigma_H=\sqrt{2/3}\,\sigma_{\mathrm{cm}}$, and the CMPL is the $2$-D horizontal Rayleigh
inflation $K_H\sigma_H$ with the representative multiplier $K_H=6.0$, the value used at
civil-aviation horizontal integrity risks in ARAIM \citep{blanch2015,wgc2016araim}. The
representative single-provider $\mathrm{HPL}_{\mathrm{ARAIM}}=\SI{15}{m}$ is an
illustrative placeholder, not a computed lunar protection level; it enters only to
exhibit the envelope's structure. The total is the triangle-inequality (arithmetic) sum
$\mathrm{HPL}_{\mathrm{ARAIM}}+\mathrm{CMPL}$ rather than a root-sum-square, because the
single-provider and cross-provider position-domain errors are not guaranteed independent
and can align (Section~\ref{sec:envelope}); the arithmetic sum is the conservative bound
and the root-sum-square would overclaim integrity. The envelope is valid only where the
cross-provider common-mode is not already absorbed into a provider's
$\sigma_{\mathrm{URA}}$ overbound (Section~\ref{sec:conditionality}); this
anti-double-counting condition is part of the model's stated scope.
